% GENERATED FILE. Edit canonical lessons, then run npm run generate:books.
% canonical-source-hash: fnv1a64:3770a5414a854b0e
% canonical-lessons: JA-C452-buchou, JA-C452-kachou, JA-C452-buka, JA-C452-shain, JA-C452-meishi

\chapter{People at Work}
\label{ch:ja-people-at-work}
\hlchaptermodality{452}

\begin{chapteropening}
\textbf{By the end of this chapter:} \emph{I can name the people I work with.}

Complete the last lesson of chapter 452: I can name the people I work with.
\end{chapteropening}

\section[\texorpdfstring{\ja{ぶちょう} (buchō)}{ぶちょう (buchō)}]{\ja{ぶちょう} (buchō) — a department head}

\label{lesson:JA-C452-buchou}



\begin{quote}
Before the new one: say the Japanese for to take a break, then the Japanese for to work overtime.
\end{quote}

\subsection*{You'll want to know: \ja{ぶちょう}}
\textbf{\ja{ぶちょう}} — \emph{buchō} — \textquotedblleft{}a department head\textquotedblright{}.

\begin{quote}
\textbf{\ja{ぶちょうは} \ja{かいぎで} \ja{はなしました。}} — \emph{buchō wa kaigi de hanashimashita} — \textquotedblleft{}The department head spoke at the meeting.\textquotedblright{}
\end{quote}

\begin{cousinweb}
\textbf{\ja{ぶちょう}} (\emph{buchō}) is Sino-Japanese (kango): two Chinese characters read with their on'yomi, readings borrowed from Chinese. They are \emph{bu}, \textquotedblleft{}section\textquotedblright{} (Middle Chinese \emph{buX}) and \emph{chō}, \textquotedblleft{}leader\textquotedblright{} (\emph{trjangX}) (KANJIDIC2; Middle Chinese in Baxter's transcription).
\end{cousinweb}

\begin{grammarlens}[title={What you've built}]
One more word for describing what you lived and felt.
\end{grammarlens}

\subsection*{Your turn}
\begin{itemize}
\raggedright
  \item \emph{Say it:} \emph{buchō}
  \item \emph{Say it:} \emph{buchō}, once more
  \item \emph{Say it:} \emph{buchō wa kaigi de hanashimashita} — The department head spoke at the meeting.
  \item \emph{Recall:} say the Japanese for to take a break, then the Japanese for to work overtime, then say \emph{buchō} again
\end{itemize}

\begin{tcolorbox}[breakable,colback=teal!4,colframe=teal!35!black,title={Before you move on}]
What does \ja{ぶちょう} mean? (\textquotedblleft{}A department head\textquotedblright{}.) Say it once more.
\end{tcolorbox}
\section[\texorpdfstring{\ja{かちょう} (kachō)}{かちょう (kachō)}]{\ja{かちょう} (kachō) — a section chief}

\label{lesson:JA-C452-kachou}



\begin{quote}
Before the new one: say the Japanese for to work overtime, then the Japanese for a department head.
\end{quote}

\subsection*{You'll want to know: \ja{かちょう}}
\textbf{\ja{かちょう}} — \emph{kachō} — \textquotedblleft{}a section chief\textquotedblright{}.

\begin{quote}
\textbf{\ja{かちょうは} \ja{まいあさ} \ja{はやく} \ja{きます。}} — \emph{kachō wa maiasa hayaku kimasu} — \textquotedblleft{}The section chief comes in early every morning.\textquotedblright{}
\end{quote}

\begin{cousinweb}
\textbf{\ja{かちょう}} (\emph{kachō}) is Sino-Japanese (kango): two Chinese characters read with their on'yomi, readings borrowed from Chinese. They are \emph{ka}, \textquotedblleft{}section\textquotedblright{} (Middle Chinese \emph{khwaH}) and \emph{chō}, \textquotedblleft{}leader\textquotedblright{} (\emph{trjangX}) (KANJIDIC2; Middle Chinese in Baxter's transcription).
\end{cousinweb}

\begin{grammarlens}[title={What you've built}]
One more word for describing what you lived and felt.
\end{grammarlens}

\subsection*{Your turn}
\begin{itemize}
\raggedright
  \item \emph{Say it:} \emph{kachō}
  \item \emph{Say it:} \emph{kachō}, once more
  \item \emph{Say it:} \emph{kachō wa maiasa hayaku kimasu} — The section chief comes in early every morning.
  \item \emph{Recall:} say the Japanese for to work overtime, then the Japanese for a department head, then say \emph{kachō} again
\end{itemize}

\begin{tcolorbox}[breakable,colback=teal!4,colframe=teal!35!black,title={Before you move on}]
What does \ja{かちょう} mean? (\textquotedblleft{}A section chief\textquotedblright{}.) Say it once more.
\end{tcolorbox}
\section[\texorpdfstring{\ja{ぶか} (buka)}{ぶか (buka)}]{\ja{ぶか} (buka) — a subordinate}

\label{lesson:JA-C452-buka}



\begin{quote}
Before the new one: say the Japanese for a department head, then the Japanese for a section chief.
\end{quote}

\subsection*{You'll want to know: \ja{ぶか}}
\textbf{\ja{ぶか}} — \emph{buka} — \textquotedblleft{}a subordinate\textquotedblright{}.

\begin{quote}
\textbf{\ja{かれの} \ja{ぶかは} \ja{よく} \ja{はたらきます。}} — \emph{kare no buka wa yoku hatarakimasu} — \textquotedblleft{}His staff work hard.\textquotedblright{}
\end{quote}

\begin{cousinweb}
\textbf{\ja{ぶか}} (\emph{buka}) is Sino-Japanese (kango): two Chinese characters read with their on'yomi, readings borrowed from Chinese. They are \emph{bu}, \textquotedblleft{}section\textquotedblright{} (Middle Chinese \emph{buX}) and \emph{ka}, \textquotedblleft{}below\textquotedblright{} (\emph{haeX}) (KANJIDIC2; Middle Chinese in Baxter's transcription).
\end{cousinweb}

\begin{grammarlens}[title={What you've built}]
One more word for describing what you lived and felt.
\end{grammarlens}

\subsection*{Your turn}
\begin{itemize}
\raggedright
  \item \emph{Say it:} \emph{buka}
  \item \emph{Say it:} \emph{buka}, once more
  \item \emph{Say it:} \emph{kare no buka wa yoku hatarakimasu} — His staff work hard.
  \item \emph{Recall:} say the Japanese for a department head, then the Japanese for a section chief, then say \emph{buka} again
\end{itemize}

\begin{tcolorbox}[breakable,colback=teal!4,colframe=teal!35!black,title={Before you move on}]
What does \ja{ぶか} mean? (\textquotedblleft{}A subordinate\textquotedblright{}.) Say it once more.
\end{tcolorbox}
\section[\texorpdfstring{\ja{しゃいん} (shain)}{しゃいん (shain)}]{\ja{しゃいん} (shain) — a company employee}

\label{lesson:JA-C452-shain}



\begin{quote}
Before the new one: say the Japanese for a section chief, then the Japanese for a subordinate.
\end{quote}

\subsection*{You'll want to know: \ja{しゃいん}}
\textbf{\ja{しゃいん}} — \emph{shain} — \textquotedblleft{}a company employee\textquotedblright{}.

\begin{quote}
\textbf{\ja{しゃいんは} \ja{みんな} \ja{まじめです。}} — \emph{shain wa minna majime desu} — \textquotedblleft{}The employees are all hard-working.\textquotedblright{}
\end{quote}

\begin{cousinweb}
\textbf{\ja{しゃいん}} (\emph{shain}) is Sino-Japanese (kango): two Chinese characters read with their on'yomi, readings borrowed from Chinese. They are \emph{sha}, \textquotedblleft{}company\textquotedblright{} (Middle Chinese \emph{dzyaeX}) and \emph{in}, \textquotedblleft{}member\textquotedblright{} (\emph{hjwen}) (KANJIDIC2; Middle Chinese in Baxter's transcription).
\end{cousinweb}

\begin{grammarlens}[title={What you've built}]
One more word for describing what you lived and felt.
\end{grammarlens}

\subsection*{Your turn}
\begin{itemize}
\raggedright
  \item \emph{Say it:} \emph{shain}
  \item \emph{Say it:} \emph{shain}, once more
  \item \emph{Say it:} \emph{shain wa minna majime desu} — The employees are all hard-working.
  \item \emph{Recall:} say the Japanese for a section chief, then the Japanese for a subordinate, then say \emph{shain} again
\end{itemize}

\begin{tcolorbox}[breakable,colback=teal!4,colframe=teal!35!black,title={Before you move on}]
What does \ja{しゃいん} mean? (\textquotedblleft{}A company employee\textquotedblright{}.) Say it once more.
\end{tcolorbox}
\section[\texorpdfstring{\ja{めいし} (meishi)}{めいし (meishi)}]{\ja{めいし} (meishi) — a business card}

\label{lesson:JA-C452-meishi}



\begin{quote}
Before the new one: say the Japanese for a subordinate, then the Japanese for a company employee.
\end{quote}

\subsection*{You'll want to know: \ja{めいし}}
\textbf{\ja{めいし}} — \emph{meishi} — \textquotedblleft{}a business card\textquotedblright{}.

\begin{quote}
\textbf{\ja{かいぎで} \ja{めいしを} \ja{わたしました。}} — \emph{kaigi de meishi o watashimashita} — \textquotedblleft{}I handed over my business card at the meeting.\textquotedblright{}
\end{quote}

\begin{cousinweb}
\textbf{\ja{めいし}} (\emph{meishi}) is Sino-Japanese (kango): two Chinese characters read with their on'yomi, readings borrowed from Chinese. They are \emph{mei}, \textquotedblleft{}name\textquotedblright{} (Middle Chinese \emph{mjieng}) and \emph{shi}, \textquotedblleft{}calling card\textquotedblright{} (\emph{tshjeH}) (KANJIDIC2; Middle Chinese in Baxter's transcription).
\end{cousinweb}

\begin{grammarlens}[title={What you've built}]
One more word for describing what you lived and felt.
\end{grammarlens}

\subsection*{Your turn}
\begin{itemize}
\raggedright
  \item \emph{Say it:} \emph{meishi}
  \item \emph{Say it:} \emph{meishi}, once more
  \item \emph{Say it:} \emph{kaigi de meishi o watashimashita} — I handed over my business card at the meeting.
  \item \emph{Recall:} say the Japanese for a subordinate, then the Japanese for a company employee, then say \emph{meishi} again
\end{itemize}

\begin{tcolorbox}[breakable,colback=teal!4,colframe=teal!35!black,title={Before you move on}]
What does \ja{めいし} mean? (\textquotedblleft{}A business card\textquotedblright{}.) Say it once more.
\end{tcolorbox}
